%%%=====================================================================%%%
%%%  frontmatter/cover.tex — Page de titre                              %%%
%%%                                                                      %%%
%%%  Le titre, le sous-titre et l'auteur viennent des macros \mbTitre,   %%%
%%%  \mbSousTitre et \mbAuteur définies en tête de ml4ebook.cls.         %%%
%%%                                                                      %%%
%%%  PARTI PRIS : ce livre est un livre de MATHÉMATIQUES, dont l'IA est  %%%
%%%  le champ d'application — pas l'inverse. La couverture montre donc   %%%
%%%  un panorama des cinq domaines traités, chacun avec sa figure ET sa  %%%
%%%  formule, l'apprentissage n'étant que le dernier des cinq :          %%%
%%%                                                                      %%%
%%%     algèbre linéaire · analyse · probabilités                        %%%
%%%              optimisation · apprentissage                            %%%
%%%                                                                      %%%
%%%  Composition en registres horizontaux :                              %%%
%%%     1. la rangée haute des trois domaines fondamentaux               %%%
%%%     2. le bandeau bleu qui porte le titre                            %%%
%%%     3. le filet ambre, le sous-titre, la promesse                    %%%
%%%     4. la rangée basse : optimisation et apprentissage               %%%
%%%     5. l'auteur                                                      %%%
%%%                                                                      %%%
%%%  Les teintes sont NETTEMENT plus soutenues que dans le corps du      %%%
%%%  livre : une couverture est vue de loin, souvent imprimée, parfois   %%%
%%%  en noir et blanc. Un trait à 30 % d'opacité y disparaît.            %%%
%%%=====================================================================%%%
\begin{titlepage}
\thispagestyle{empty}

% Styles locaux : ils ne servent qu'ici, mais évitent de répéter dix fois
% les mêmes réglages de police pour les formules et les intitulés.
\tikzset{
  cov formule/.style={font=\small, text=mbInk},
  cov domaine/.style={font=\sffamily\scriptsize\bfseries, text=mbSlate},
  cov courbe/.style={draw=mbBlue, line width=1.1pt},
  cov axe/.style={draw=mbSlate!75, line width=0.6pt,
                  -{Stealth[round,length=1.7mm,width=1.3mm]}},
}

\begin{tikzpicture}[remember picture,overlay,mbfig]

  %--------------------------------------------------------------------
  % Fond
  %--------------------------------------------------------------------
  \fill[mbPaper] (current page.south west) rectangle (current page.north east);

  %====================================================================
  % 1. RANGÉE HAUTE — les trois piliers mathématiques
  %====================================================================
  \coordinate (rhaut) at ($(current page.north)+(0cm,-4.6cm)$);

  %--------------------------------------------------------------------
  % 1a. ALGÈBRE LINÉAIRE — une matrice, avec ses coefficients et ses
  %     crochets. Les cases colorées la rendent lisible de loin, les
  %     nombres disent que c'est bien un objet mathématique.
  %--------------------------------------------------------------------
  \begin{scope}[shift={($(rhaut)+(-4.5cm,0cm)$)}]
    \foreach \i in {0,...,3}{
      \foreach \j in {0,...,3}{
        \pgfmathtruncatemacro{\val}{mod(3*\i+5*\j+7,9)-4}
        \ifnum\i=\j
          \fill[mbAmber!45] (-1.16+\j*0.58, 1.16-\i*0.58) rectangle ++(0.58,-0.58);
        \else
          \fill[mbBlue!\the\numexpr8+4*\i+5*\j\relax]
                (-1.16+\j*0.58, 1.16-\i*0.58) rectangle ++(0.58,-0.58);
        \fi
        \draw[mbBlue!55, line width=0.4pt]
              (-1.16+\j*0.58, 1.16-\i*0.58) rectangle ++(0.58,-0.58);
        \node[font=\scriptsize, text=mbInk]
              at (-0.87+\j*0.58, 0.87-\i*0.58) {$\val$};
      }
    }
    % Crochets de matrice
    \draw[mbInk, line width=1pt]
      (-1.32,1.16) -- (-1.32,-1.16) (-1.32,1.16) -- (-1.10,1.16)
      (-1.32,-1.16) -- (-1.10,-1.16);
    \draw[mbInk, line width=1pt]
      (1.32,1.16) -- (1.32,-1.16) (1.32,1.16) -- (1.10,1.16)
      (1.32,-1.16) -- (1.10,-1.16);
    \node[cov formule] at (0,-1.78) {$A\vect{x} = \vect{b}$};
    \node[cov domaine] at (0,-2.28) {ALGÈBRE LINÉAIRE};
  \end{scope}

  %--------------------------------------------------------------------
  % 1b. ANALYSE — une courbe et sa tangente au sommet. La tangente y est
  %     horizontale : c'est exactement la condition f'(x) = 0, celle qui
  %     fonde toute l'optimisation. La figure annonce la suivante.
  %--------------------------------------------------------------------
  \begin{scope}[shift={($(rhaut)+(0cm,0cm)$)}]
    \draw[cov axe] (-1.98,-0.78) -- (2.05,-0.78);
    \draw[cov axe] (-1.8125,-0.95) -- (-1.8125,1.02);
    \draw[cov courbe, smooth, samples=80, domain=0:2.9]
      plot ({\x*1.25-1.8125}, {(0.75+0.5*sin(deg(1.9*\x)))*1.25-0.78});
    % La tangente au maximum : elle est HORIZONTALE, c'est-à-dire f'(x)=0
    \draw[draw=mbAmber!85!black, line width=1.2pt] (-1.5,0.7825) -- (-0.05,0.7825);
    \draw[mbSlate!55, line width=0.5pt, dashed] (-0.78,0.7825) -- (-0.78,-0.78);
    \fill[mbInk] (-0.78,0.7825) circle (2pt);
    \node[cov formule] at (0,-1.78) {$f'(x) = 0$};
    \node[cov domaine] at (0,-2.28) {ANALYSE};
  \end{scope}

  %--------------------------------------------------------------------
  % 1c. PROBABILITÉS — la courbe en cloche, avec l'aire sous la courbe :
  %     l'image même d'une probabilité.
  %--------------------------------------------------------------------
  \begin{scope}[shift={($(rhaut)+(4.5cm,0cm)$)}]
    % L'aire grisée court sur un écart-type de part et d'autre de la
    % moyenne : c'est une probabilité, montrée comme une surface.
    \fill[mbBlue!35]
      (-0.7375,-0.78) --
      plot[smooth, samples=50, domain=0.91:2.09]
        ({\x*1.25-1.875}, {1.25*exp(-(\x-1.5)^2/0.7)*1.25-0.78}) --
      (0.7375,-0.78) -- cycle;
    \draw[cov axe] (-2.05,-0.78) -- (2.1,-0.78);
    \draw[cov courbe, smooth, samples=80, domain=0:3]
      plot ({\x*1.25-1.875}, {1.25*exp(-(\x-1.5)^2/0.7)*1.25-0.78});
    \draw[mbAmber!85!black, line width=0.9pt, dashed] (0,-0.78) -- (0,0.78);
    \node[font=\scriptsize, text=mbAmber!80!black, below] at (0,-0.81) {$\mu$};
    \node[cov formule] at (0,-1.78) {$\Normale(\mu,\sigma^{2})$};
    \node[cov domaine] at (0,-2.28) {PROBABILITÉS};
  \end{scope}

  %====================================================================
  % 2. Bandeau du titre
  %    ATTENTION : les deux coins doivent avoir des ordonnées DIFFÉRENTES.
  %    La version d'origine allait de `current page.west` à
  %    `current page.east|-current page.center` — or `west` est déjà à la
  %    hauteur du centre : le rectangle était plat, donc invisible, et le
  %    titre blanc se retrouvait posé sur le fond crème, illisible.
  %====================================================================
  \fill[mbBlue] (current page.west|-current page.center)
        rectangle ([yshift=3.7cm]current page.east|-current page.center);
  \fill[mbAmber] (current page.west|-current page.center)
        rectangle ([yshift=-0.28cm]current page.east|-current page.center);

  \node[anchor=west,text width=11.5cm,align=left]
        at ([yshift=1.85cm,xshift=1.9cm]current page.west) {%
    \sffamily\bfseries\color{white}\fontsize{34}{39}\selectfont
    Machine Learning\\ For Engineer};

  %====================================================================
  % 3. Sous-titre et promesse, sous le filet — donc sur le fond crème
  %====================================================================
  \node[anchor=north west,text width=12.4cm,align=left]
        at ([yshift=-1.05cm,xshift=1.9cm]current page.west) {%
    \sffamily\color{mbInk}\large \mbSousTitre};
  \node[anchor=north west,text width=12.4cm,align=left]
        at ([yshift=-2.15cm,xshift=1.9cm]current page.west) {%
    \sffamily\bfseries\color{mbAmber!80!black}\normalsize
    Aucun prérequis : on part vraiment de zéro.};

  %====================================================================
  % 4. RANGÉE BASSE — ce qu'on construit avec les trois piliers
  %====================================================================
  \coordinate (rbas) at ($(current page.south)+(0cm,6.5cm)$);

  %--------------------------------------------------------------------
  % 4a. OPTIMISATION — lignes de niveau d'une fonction de coût et la
  %     descente qui rejoint le minimum. Le zigzag n'est pas décoratif :
  %     c'est le comportement réel d'une descente dans une vallée
  %     allongée, l'amplitude décroissant à chaque pas.
  %--------------------------------------------------------------------
  \begin{scope}[shift={($(rbas)+(-3.4cm,0cm)$)}]
    \begin{scope}[rotate=20]
      \foreach \r/\t in {1.15/34, 0.92/44, 0.69/54, 0.46/64, 0.23/74}{
        \draw[mbBlue!\t, line width=0.8pt]
              (0,0) ellipse ({\r*1.9cm} and {\r*1.08cm});
      }
      \fill[mbBlue] (0,0) circle (2pt);
      \draw[-{Stealth[round,length=2.2mm,width=1.7mm]},
            draw=mbAmber!85!black, line width=1.1pt]
        (-1.95,0.55) -- (-1.03,-0.46) -- (-0.55,0.32)
        -- (-0.27,-0.18) -- (-0.11,0.10) -- (-0.02,-0.01);
      \foreach \p in {(-1.95,0.55), (-1.03,-0.46), (-0.55,0.32),
                      (-0.27,-0.18), (-0.11,0.10)}{
        \fill[mbAmber!85!black] \p circle (1.8pt);
      }
    \end{scope}
    \node[cov formule] at (0,-1.85)
         {$\params \leftarrow \params - \pas\,\grad\cout(\params)$};
    \node[cov domaine] at (0,-2.35) {OPTIMISATION};
  \end{scope}

  %--------------------------------------------------------------------
  % 4b. APPRENTISSAGE — le réseau de neurones. Il vient en dernier, et
  %     il n'est qu'une application de tout ce qui précède : une
  %     multiplication par une matrice, puis une fonction, le tout réglé
  %     par une descente de gradient.
  %--------------------------------------------------------------------
  % `minimum size` doit être posé APRÈS le style `neurone`, dans la liste
  % d'options du nœud : passé par `every node`, il serait écrasé par le
  % 7 mm du style — plus large que l'espacement, et les neurones se
  % chevaucheraient.
  \begin{scope}[shift={($(rbas)+(3.4cm,0cm)$)}]
    \foreach \i in {0,1,2}
      \node[neurone entree, draw=mbAzure, fill=mbAzure!25, line width=0.8pt,
            minimum size=4.4mm]
            (re\i) at (-1.5, {0.62*(1-\i)}) {};
    \foreach \i in {0,1,2,3}
      \node[neurone cache, draw=mbBlue, fill=mbBlue!20, line width=0.8pt,
            minimum size=4.4mm]
            (rc\i) at (0, {0.62*(1.5-\i)}) {};
    \foreach \i in {0,1}
      \node[neurone sortie, draw=mbAmber!85!black, fill=mbAmber!45,
            line width=0.8pt, minimum size=4.4mm]
            (rs\i) at (1.5, {0.62*(0.5-\i)}) {};
    \foreach \a in {0,1,2} \foreach \b in {0,1,2,3}
      \draw[mbSlate!45, line width=0.45pt] (re\a) -- (rc\b);
    \foreach \a in {0,1,2,3} \foreach \b in {0,1}
      \draw[mbSlate!45, line width=0.45pt] (rc\a) -- (rs\b);
    \node[cov formule] at (0,-1.85)
         {$\vect{y} = \sigma(W\vect{x} + \vect{b})$};
    \node[cov domaine] at (0,-2.35) {APPRENTISSAGE};
  \end{scope}

  %====================================================================
  % 5. Auteur
  %====================================================================
  \node[anchor=south west,text width=12.4cm,align=left]
        at ([yshift=1.5cm,xshift=1.9cm]current page.south west) {%
    \sffamily
    {\bfseries\color{mbInk}\large\mbAuteur}\\[6pt]
    {\color{mbBlue}\footnotesize\faIcon{calendar-alt}\ \the\year}};
\end{tikzpicture}
\end{titlepage}

% Verso de la page de titre : colophon
\thispagestyle{empty}
\null\vfill
{\footnotesize\sffamily\color{mbSlate}
\noindent Édition \the\year. Document composé avec \LaTeX{} (pdf\TeX), polices
Palatino (texte et mathématiques), TeX~Gyre~Heros (titres) et DejaVu~Sans~Mono
(code).
Figures et schémas réalisés avec \textsf{TikZ} et \textsf{pgfplots}, coloration
syntaxique assurée par \textsf{minted}/Pygments, encadrés par
\textsf{tcolorbox}, bibliographie par \textsf{biblatex}.\par
\medskip
\noindent Ce livre rassemble en un seul volume les mathématiques dont
l'intelligence artificielle et la science des données ont réellement besoin :
algèbre linéaire, analyse, probabilités, statistique et optimisation. Il est
écrit pour être lu sans aucun prérequis : on part de zéro, on explique avec
des mots avant d'écrire des symboles, et on n'avance jamais sans avoir
d'abord dit pourquoi.\par}
\clearpage
